% problem-000324 7 超四面体连接结构
\section*{7. 超四面体连接结构}
T_{2}
（图：）

（图：）

（图：）

（图：）

（图：）

（图：）

\subsection*{7-1}
上图中 $\mathbf { T } _ { 1 } .$ 、 $\mathbf { T } _ { 2 }$ 和 $\mathbf { T } _ { 3 }$ 的化学式分别为 $\mathbf { A X } _ { 4 }$ $\mathbf { A } _ { 4 } \mathbf { X } _ { 1 0 }$ 和 $\mathbf { A } _ { 1 0 } \mathbf { X } _ { 2 0 } ,$ ，推出超四⾯体 $\mathrm { T } _ { 4 }$ 的化学式。

\subsection*{7-2}
分别指出超四⾯体 $\mathrm { { T } _ { 3 } } .$ 、 $\mathrm { T } _ { 4 }$ 中各有⼏种环境不同的X原⼦，每种X原⼦各连接⼏个A原⼦？在上述两种超四⾯体中每种X原⼦的数⽬各是多少？

\subsection*{7-3}
若分别以 $\mathbf { T } _ { 1 \mathrm { : } }$ $\mathbf { T } _ { 2 } ,$ $\mathrm { { T } _ { 3 \mathrm { { \cdot } } } }$ $\mathrm { { T } _ { 4 } }$ 为结构单元共顶点相连（顶点X原⼦只连接两个A原⼦），形成⽆限三维结构，分别写出所得三维⻣架的化学式。

\subsection*{7-4}
欲使上述 $\mathbf { T } _ { 3 }$ 超四⾯体连接所得三维⻣架的化学式所带电荷分别为+4、0和−4，A选 $\mathrm { Z n ^ { 2 + } }$ 、 $\mathrm { I n ^ { 3 + } }$ 或 $\mathrm { G e ^ { 4 + } }$ X 取 $S ^ { 2 - }$ ，给出带三种不同电荷的⻣架的化学式（各给出⼀种，结构单元中的离⼦数成简单整数⽐）。
